\section{Conclusion}\label{sec:conclusion}
Neural Bellman Operators provide an evaluation and improvement structure for continuous-time economic control and equilibrium problems. The critic approximates the continuation value associated with a specified policy; the actor makes feasible choices using that value; and independent verification measures the remaining economic error. Differential verification, monotone Bellman approximation, and complete-action enclosures make these roles explicit across the models considered here.

The capital application connects a fitted policy to the original diffusion through a policy-specific payoff comparison. The analytical schedule supplies an absolute benchmark, and direct paired inference measures the gain from a particular learned policy after accounting for implementation and sampling error. Costate accuracy provides a mathematical link between critic evaluation and Hamiltonian improvement. The computational comparisons assess that link against direct policy optimization and other feasible alternatives, while reporting the work required both to construct and to verify each candidate.

The economic program remains the one that motivates NBO. Endogenous preferences combine consumption, portfolio risk, and costly internal adjustment. Recursive utility changes the continuation-value equation. Temporal selves and competing firms require equilibrium-specific deviation conditions. Their applications retain their own models, derivations, economic comparisons, and verification results. The resulting framework makes quantitative claims assessable in utility, compensation, and unilateral-profit units, with the information structure and approximation domain stated for each claim.

The theoretical and computational results identify where further precision is valuable: policy evaluation must improve the derivatives relevant for decisions, and the resulting gain must justify the total work and implementation cost. This connection between a numerical Bellman calculation and a verified economic decision is the organizing contribution of the paper.
